\documentclass[11pt,a4paper]{article}

% ---------- Encoding e idioma ----------
\usepackage[utf8]{inputenc}
\usepackage[T1]{fontenc}
\usepackage[spanish,es-tabla]{babel}

% ---------- Diseño de página ----------
\usepackage[margin=2.5cm]{geometry}
\usepackage{setspace}
\onehalfspacing

% ---------- Matemáticas ----------
\usepackage{amsmath,amssymb,amsthm,mathtools}

% ---------- Figuras y tablas ----------
\usepackage{graphicx}
\usepackage{booktabs}
\usepackage{caption}
\usepackage{float}

% ---------- Color ----------
\usepackage{xcolor}
\definecolor{azulPSNR}{rgb}{0.10,0.35,0.65}
\definecolor{verdeSSIM}{rgb}{0.10,0.55,0.25}

% ---------- Enlaces ----------
\usepackage[colorlinks=true, linkcolor=blue!60!black, urlcolor=blue!60!black]{hyperref}

% ---------- Metadatos ----------
\title{
    \vspace{-1.5cm}
    \textbf{Métricas de Calidad de Imagen} \\[0.3cm]
    \Large PSNR y SSIM: Definición, Fórmulas e Interpretación
}
\author{Jamin Yauri Cajas}
\date{\today}

% ============================================================
\begin{document}
\maketitle

\begin{abstract}
\noindent
Este documento explica de forma clara las dos métricas más usadas en procesamiento
de imágenes para comparar una imagen procesada contra una de referencia:
\textbf{PSNR} (Peak Signal-to-Noise Ratio) y \textbf{SSIM} (Structural Similarity
Index Measure). Se presentan sus fórmulas matemáticas, el significado intuitivo
de cada término, ejemplos numéricos, sus rangos típicos y, sobre todo, las
limitaciones que cada una arrastra. Finalmente se interpreta su significado
particular en el contexto del pipeline de preprocesamiento de imágenes de
escritorio, donde ambas métricas se aplican entre la imagen \emph{original} y
la imagen \emph{preprocesada}.
\end{abstract}

% ============================================================
\section{¿Qué problema resuelven estas métricas?}

Cuando se modifica una imagen — por compresión, por ruido, o por un pipeline de
preprocesamiento como el nuestro — necesitamos una forma \textbf{numérica} de
responder a la pregunta: \emph{¿qué tan parecida es la imagen modificada a la
original?} Mirar con los ojos funciona para una imagen; para evaluar miles
necesitamos métricas objetivas.

PSNR y SSIM son las dos métricas clásicas, y son \textbf{complementarias}:

\begin{itemize}
    \item \textbf{PSNR} mide la diferencia \emph{píxel a píxel}. Es rápida, intuitiva,
          pero ciega a cómo percibe el ojo humano.
    \item \textbf{SSIM} mide la similitud \emph{estructural}. Está inspirada en la
          visión humana: si dos imágenes tienen la misma estructura (bordes,
          contraste, texturas) aunque los píxeles difieran, SSIM las considera
          parecidas.
\end{itemize}

% ============================================================
\section{\textcolor{azulPSNR}{PSNR} — Peak Signal-to-Noise Ratio}

\subsection{Definición intuitiva}

PSNR responde la pregunta: \emph{¿qué tan alta es la señal comparada con el
error introducido?} Se expresa en \textbf{decibelios (dB)}, una escala
logarítmica: más dB = menos diferencia = mejor.

\subsection{Fórmula matemática}

PSNR se calcula a partir del \textbf{error cuadrático medio} (MSE). Dadas dos
imágenes \(I\) y \(K\) de tamaño \(H \times W\):

\begin{equation}
    \operatorname{MSE}(I, K) = \frac{1}{HW} \sum_{i=1}^{H} \sum_{j=1}^{W}
    \bigl[ I(i,j) - K(i,j) \bigr]^2.
    \label{eq:mse}
\end{equation}

\vspace{0.3cm}
\noindent
El MSE es simplemente el \textbf{promedio de las diferencias al cuadrado} entre
cada par de píxeles. Si las dos imágenes son idénticas, MSE = 0. Si son muy
distintas, el MSE crece.

A partir del MSE se calcula PSNR:

\begin{equation}
    \boxed{\;
    \operatorname{PSNR}(I, K) = 10 \cdot \log_{10}\!\left(\frac{\operatorname{MAX}^2}{\operatorname{MSE}(I,K)}\right)
    \;}
    \label{eq:psnr}
\end{equation}

donde \(\operatorname{MAX}\) es el valor máximo que puede tomar un píxel:
\(\operatorname{MAX} = 255\) para imágenes \texttt{uint8}, o \(\operatorname{MAX} = 1\)
para imágenes \texttt{double} en \([0,1]\) (como en nuestro pipeline).

\subsection{¿Por qué logarítmica?}

Porque la percepción humana del ruido también lo es: la diferencia entre una
imagen con MSE = 0.01 y una con MSE = 0.001 es enorme a la vista, y el logaritmo
la amplifica de forma legible. Además, expresar la relación en decibelios
permite comparar órdenes de magnitud muy distintas en una escala cómoda.

\subsection{Ejemplo numérico}

Supongamos dos imágenes de \(2 \times 2\) píxeles, en rango \([0,1]\):

\[
I = \begin{pmatrix} 0.5 & 0.4 \\ 0.6 & 0.5 \end{pmatrix}, \qquad
K = \begin{pmatrix} 0.52 & 0.39 \\ 0.61 & 0.48 \end{pmatrix}.
\]

Cálculo del MSE:
\begin{align*}
\operatorname{MSE} &= \frac{1}{4}\bigl[(0.02)^2 + (0.01)^2 + (0.01)^2 + (0.02)^2\bigr] \\
                   &= \frac{1}{4}\bigl[0.0004 + 0.0001 + 0.0001 + 0.0004\bigr] = 0.00025.
\end{align*}

PSNR con \(\operatorname{MAX} = 1\):
\[
\operatorname{PSNR} = 10 \cdot \log_{10}\!\left(\frac{1}{0.00025}\right) = 10 \cdot \log_{10}(4000) \approx 36{,}02 \text{ dB}.
\]

\subsection{Rangos típicos e interpretación}

\begin{table}[H]
\centering
\caption{Interpretación cualitativa de PSNR (imágenes naturales).}
\begin{tabular}{@{}cl@{}}
\toprule
PSNR (dB) & Calidad percibida \\
\midrule
\(> 40\)      & Virtualmente idéntica (compresión de alta calidad) \\
\(30 \text{ a } 40\) & Muy similar (pocas diferencias visibles) \\
\(20 \text{ a } 30\) & Diferencias notables (pipeline activo) \\
\(< 20\)      & Diferencias muy fuertes (fuerte modificación o error) \\
\bottomrule
\end{tabular}
\end{table}

\subsection{Limitaciones importantes}

PSNR tiene dos debilidades críticas:

\begin{enumerate}
    \item \textbf{Es ciego a la estructura.} Si a una foto le desplazas todos los
          píxeles un solo píxel hacia la derecha, visualmente sigue siendo
          prácticamente la misma imagen, pero el MSE es enorme y el PSNR se
          desploma.
    \item \textbf{No modela el ojo humano.} Las personas percibimos peor el error
          en regiones texturizadas que en regiones planas; PSNR trata ambas por
          igual.
\end{enumerate}

Por eso, en investigación moderna se combina PSNR con SSIM o con métricas
perceptuales basadas en redes neuronales (LPIPS, DISTS).

% ============================================================
\section{\textcolor{verdeSSIM}{SSIM} — Structural Similarity Index Measure}

\subsection{Definición intuitiva}

SSIM responde una pregunta muy distinta a PSNR: \emph{¿las dos imágenes tienen
la misma estructura?} Está inspirada en el \textbf{Sistema Visual Humano}, que
no mira píxel por píxel, sino que percibe:

\begin{itemize}
    \item La \textbf{luminancia} (brillo promedio).
    \item El \textbf{contraste} (cuánto varía la intensidad).
    \item La \textbf{estructura} (los patrones y bordes).
\end{itemize}

SSIM produce un valor entre \(-1\) y \(+1\):
\begin{itemize}
    \item \(\operatorname{SSIM} = 1\) : imágenes idénticas.
    \item \(\operatorname{SSIM} \approx 0\) : sin relación estructural.
    \item \(\operatorname{SSIM} < 0\) : imágenes estructuralmente opuestas (raro).
\end{itemize}

En la práctica, valores por encima de \(0{,}9\) indican imágenes visualmente
equivalentes.

\subsection{Fórmula matemática}

SSIM se calcula sobre ventanas (parches) de la imagen. Para dos parches \(x\) y
\(y\) extraídos en la misma posición de las dos imágenes:

\begin{equation}
    \boxed{\;
    \operatorname{SSIM}(x,y) =
    \underbrace{\left[\frac{2\mu_x \mu_y + c_1}{\mu_x^2 + \mu_y^2 + c_1}\right]}_{l(x,y)\,:\text{ luminancia}} \cdot
    \underbrace{\left[\frac{2\sigma_x \sigma_y + c_2}{\sigma_x^2 + \sigma_y^2 + c_2}\right]}_{c(x,y)\,:\text{ contraste}} \cdot
    \underbrace{\left[\frac{\sigma_{xy} + c_3}{\sigma_x \sigma_y + c_3}\right]}_{s(x,y)\,:\text{ estructura}}
    \;}
    \label{eq:ssim}
\end{equation}

donde:
\begin{itemize}
    \item \(\mu_x, \mu_y\): medias de cada parche (\textbf{luminancia promedio}).
    \item \(\sigma_x, \sigma_y\): desviaciones estándar (\textbf{contraste local}).
    \item \(\sigma_{xy}\): covarianza entre los parches (\textbf{patrón estructural compartido}).
    \item \(c_1, c_2, c_3\): constantes pequeñas que estabilizan la fórmula cuando los denominadores son casi cero.
\end{itemize}

El SSIM global de la imagen es el \textbf{promedio} de los SSIM de todas las
ventanas.

\subsection{Interpretación de los tres componentes}

La fórmula se lee como el producto de tres comparaciones \textbf{independientes}:

\begin{enumerate}
    \item \(l(x,y)\) compara \textbf{brillos}: si \(\mu_x \approx \mu_y\), este
          término vale \(\approx 1\).
    \item \(c(x,y)\) compara \textbf{contrastes}: si \(\sigma_x \approx \sigma_y\),
          este término vale \(\approx 1\).
    \item \(s(x,y)\) compara \textbf{estructuras}: si los patrones locales
          coinciden (mismo borde, misma textura), este término vale \(\approx 1\).
\end{enumerate}

El producto final de los tres es 1 solo si los tres aspectos coinciden.
\textbf{Este es el poder de SSIM}: separa brillo, contraste y estructura, y un
cambio en cualquiera de ellos se penaliza de forma independiente y explícita.

\subsection{Ejemplo conceptual: por qué SSIM es mejor que PSNR}

Imagina dos modificaciones a una foto:

\begin{enumerate}
    \item \textbf{Suma un valor constante} a todos los píxeles (subir el brillo
          \(+0{,}1\)).
    \item \textbf{Introduce ruido aleatorio} con amplitud \(\pm 0{,}1\).
\end{enumerate}

Ambas producen el \textbf{mismo MSE} y, por tanto, el \textbf{mismo PSNR}. Pero:

\begin{itemize}
    \item El primer caso (brillo subido) \textbf{conserva la estructura}: solo
          ves la foto más clara. El SSIM es alto.
    \item El segundo caso (ruido) \textbf{destruye la estructura}. El SSIM cae.
\end{itemize}

Esto coincide con la percepción humana: el ruido molesta, una corrección de
brillo no.

\subsection{Rangos típicos e interpretación}

\begin{table}[H]
\centering
\caption{Interpretación cualitativa de SSIM.}
\begin{tabular}{@{}cl@{}}
\toprule
SSIM & Calidad percibida \\
\midrule
\(> 0{,}95\)        & Visualmente idénticas \\
\(0{,}85 - 0{,}95\) & Diferencias sutiles, estructura intacta \\
\(0{,}70 - 0{,}85\) & Diferencias notables pero reconocible \\
\(< 0{,}70\)        & Cambios estructurales fuertes \\
\bottomrule
\end{tabular}
\end{table}

\subsection{Limitaciones de SSIM}

\begin{enumerate}
    \item \textbf{Sensible al tamaño de ventana.} El valor depende del tamaño del
          parche utilizado (típicamente \(11\times 11\) gaussiano).
    \item \textbf{Trabaja en escala de grises.} Para imágenes a color, se aplica
          sobre la luminancia. Un cambio de tinte sin cambio de estructura no lo
          detecta bien.
    \item \textbf{No es distancia.} No se puede interpretar como porcentaje
          riguroso; es un índice de similitud perceptual.
\end{enumerate}

% ============================================================
\section{PSNR vs. SSIM — tabla comparativa}

\begin{table}[H]
\centering
\caption{Comparación directa de PSNR y SSIM.}
\renewcommand{\arraystretch}{1.3}
\begin{tabular}{@{}p{0.28\textwidth} p{0.32\textwidth} p{0.32\textwidth}@{}}
\toprule
 & \textbf{PSNR} & \textbf{SSIM} \\
\midrule
\textbf{Qué mide} & Diferencia píxel a píxel & Similitud estructural \\
\textbf{Unidad} & Decibelios (dB) & Índice sin unidad \\
\textbf{Rango} & \([0, \infty)\); ideal \(> 30\) & \([-1, 1]\); ideal \(> 0{,}95\) \\
\textbf{Interpretación} & Cuantitativa, dB & Cualitativa, perceptual \\
\textbf{Sensibilidad humana} & Baja & Alta (diseñada para el ojo) \\
\textbf{Cálculo} & Simple, rápido & Requiere ventanas deslizantes \\
\textbf{Detecta cambios de brillo} & Sí (penaliza) & Poco (los tolera) \\
\textbf{Detecta ruido} & Sí & Sí (más fuerte) \\
\textbf{Detecta desplazamiento} & Muy penalizante & Menos penalizante \\
\bottomrule
\end{tabular}
\end{table}

% ============================================================
\section{Interpretación en este proyecto}

En el pipeline de preprocesamiento, las dos métricas se calculan entre la
imagen \textbf{original} y la imagen \textbf{preprocesada}. No existe un
\emph{ground truth} (imagen perfecta) con el cual compararlas.

Esto implica una interpretación \textbf{invertida} respecto al uso clásico:

\begin{itemize}
    \item En un contexto de compresión o denoising, se quiere un PSNR/SSIM
          \textbf{alto}: la imagen reconstruida debe parecerse lo más posible
          a la original limpia.
    \item En nuestro contexto, un PSNR/SSIM \textbf{bajo} significa que el
          pipeline \textbf{modificó mucho} la imagen de entrada. Esto es bueno
          cuando la imagen original tenía defectos (ruido, mal balance de blancos,
          iluminación irregular): el pipeline está trabajando.
\end{itemize}

\subsection{Ejemplos concretos del dataset}

De la Tabla de métricas del informe principal:

\begin{itemize}
    \item \texttt{p06\_stationery2}: PSNR 29.40, SSIM 0.964. Imagen con buena
          iluminación y color balanceado. El pipeline apenas la modifica porque
          ya estaba bien.
    \item \texttt{p08\_stationery4}: PSNR 13.66, SSIM 0.599. Imagen con fuerte
          dominante cromática y mala iluminación. El pipeline la modifica
          agresivamente, por eso ambas métricas son bajas — no es un fallo,
          es el comportamiento esperado.
    \item \texttt{p12\_markers\_holder}: PSNR 22.88, SSIM 0.353. Aquí el PSNR
          es moderado pero el SSIM es bajo — el pipeline introdujo cambios
          estructurales (probablemente la ecualización reveló texturas antes
          sumidas en sombra).
\end{itemize}

\subsection{¿Qué métrica se debe creer más?}

Para este proyecto, \textbf{SSIM es la métrica más informativa} porque:

\begin{enumerate}
    \item Es más sensible a cambios estructurales reales (las cosas que importan
          para la detección posterior).
    \item Correlaciona mejor con la percepción humana de ``¿se ve mejor?''.
    \item Distingue cambios de brillo (benéficos) de cambios de estructura
          (potencialmente dañinos).
\end{enumerate}

PSNR se reporta porque es estándar y da una referencia dimensional, pero en un
pipeline de \textbf{realce} (no de reconstrucción) su valor absoluto es menos
crítico que el de SSIM.

% ============================================================
\section{Conclusión}

PSNR y SSIM miden dos aspectos distintos de la similitud entre imágenes:
\textbf{PSNR} mide la diferencia cuantitativa píxel a píxel en decibelios,
ideal por su simplicidad y rapidez pero ciega a la estructura; \textbf{SSIM}
mide la similitud perceptual descomponiendo la comparación en tres componentes
(luminancia, contraste, estructura), más costosa pero mucho más coherente con
la visión humana. Ambas se reportan juntas porque son complementarias: PSNR
da la magnitud, SSIM da el significado. En el contexto del preprocesamiento
de imágenes de escritorio, su interpretación se invierte respecto al uso
clásico — valores bajos no son un defecto, sino evidencia de que el algoritmo
está corrigiendo defectos reales de la imagen de entrada.

\end{document}
